\chapter{Breast Milk Storage}

\section*{Definition and Overview}
This chapter covers the safe storage of expressed breast milk -- how to collect, store, thaw, and use breast milk so that it retains its nutritional and protective properties and so that the risk of bacterial contamination is kept as low as possible. It concerns not a disease but the practical care of a food that is given to a baby. Expressed milk is commonly stored when a mother returns to work or study, when she is separated from her baby, or when she is building a supply for later use, and many mothers express and freeze milk routinely. Following simple, consistent hygiene and temperature guidelines prevents waste and, in rare cases, illness in the baby, and allows a mother to provide breast milk even when she and her baby are apart.

\section*{Epidemiology}
Most breastfeeding women express milk at some stage, and many store it regularly, whether occasionally or as their principal method of feeding. Safe handling and storage matter because expressed milk is food for a young baby whose immune system is immature, and breast milk is a biological fluid that can support bacterial growth if mishandled. With good hygiene and correct temperature control, contamination is very unusual. The main reasons milk is wasted are uncertainty about how long it can be kept, thawing more than is needed, and confusion about normal changes in the milk's appearance or smell, which are frequently mistaken for spoilage.

\section*{Aetiology and Pathophysiology}
Breast milk contains nutrients, antibodies, enzymes, and living cells that help protect and nourish the baby. These components remain stable when chilled or frozen, but bacteria from the mother's skin, the pump, the hands, or the containers can contaminate milk, and warmth encourages their growth. Storing milk at the correct temperature slows or stops that growth: refrigeration markedly slows bacterial multiplication, and freezing effectively halts it. Some nutritional and immune properties are gradually lost even in the freezer, so it is recommended to use the oldest stored milk first. Milk naturally separates into a cream layer and a watery layer on standing, and its colour can vary with the mother's diet and the stage of the feed -- these are normal features and do not mean the milk is spoiled.

\section*{Clinical Features}
There are no symptoms of disease, but it is useful to recognise how good milk should look, smell, and behave:

\begin{itemize}
    \item Fresh milk may appear bluish, yellowish, or creamy, and the colour varies between expressions and feeds
    \item Stored milk separates into layers; gentle shaking mixes them back together
    \item Some mothers' milk has a soapy smell caused by the enzyme lipase; it remains perfectly safe to feed
    \item Spoiled milk smells sour or rancid and may appear curdled or clotted
    \item Milk with an unusual appearance, smell, or taste should not be used
\end{itemize}

\section*{Differential Diagnosis}
\begin{itemize}
    \item Normal variation in milk colour or smell versus true spoilage
    \item Separated fat and watery layers versus curdled, spoiled milk
    \item A soapy smell from lipase activity versus bacterial contamination
    \item Small amounts of fat adhering to the container versus lumps of curdled milk
\end{itemize}

\section*{When to Seek Medical Care}
See a doctor, midwife, or maternal and child health nurse if your baby becomes unwell after feeding expressed milk -- especially with fever, vomiting, diarrhoea, or poor feeding -- or if you are unsure whether stored milk is safe to use. It is always better to seek advice and discard milk than to take a risk with a young baby.

\textbf{Call 000 (or local emergency services) for:}
\begin{itemize}
    \item A baby who is floppy, lethargic, or difficult to wake
    \item A baby who is not feeding at all or shows signs of dehydration, such as a sunken fontanelle or very few wet nappies
    \item Repeated forceful vomiting or a high fever in a young baby
    \item Difficulty breathing or blue-tinged lips
\end{itemize}

\section*{Diagnosis and Investigations}
There are no laboratory tests required for routine home storage. Instead, assessment means following the storage guidelines and using simple checks:

\begin{itemize}
    \item \textbf{Label and date:} always mark each container with the date the milk was expressed
    \item \textbf{Room temperature (under about 25 degrees Celsius):} up to 4 hours is commonly advised
    \item \textbf{Refrigerator (0--4 degrees Celsius):} around 3 days is widely recommended, although some sources allow up to 5--8 days
    \item \textbf{Freezer (about minus 18 degrees Celsius):} 6 months is often given as the optimum, with up to 12 months acceptable; milk stored longer is not unsafe but may lose some nutritional value
    \item \textbf{Thawing:} in the refrigerator or in a bowl of warm water; a microwave must never be used
    \item \textbf{After a feed:} discard any milk left in the bottle after the baby has fed from it
\end{itemize}

Guidelines vary slightly between organisations, so follow the advice of your local health service or maternity unit, which will be consistent and current.

\section*{Management}
\begin{itemize}
    \item \textbf{Wash hands} thoroughly before expressing or handling milk
    \item \textbf{Clean equipment:} wash and sterilise pump parts and containers after each use, or wash them in hot, soapy water
    \item \textbf{Containers:} use clean, food-grade, BPA-free bottles or breast milk storage bags designed for freezing
    \item \textbf{Store in the right place:} place milk at the back of the fridge or freezer, not in the door, where the temperature fluctuates most
    \item \textbf{Leave space:} milk expands when frozen, so do not overfill containers
    \item \textbf{Cool quickly:} refrigerate freshly expressed milk promptly
    \item \textbf{Use the oldest milk first} and rotate the stock
    \item \textbf{Do not refreeze:} once thawed, milk should be used within 24 hours and must not be frozen again
    \item \textbf{Transport:} carry stored milk in a cool bag or insulated container with an ice brick when travelling
\end{itemize}

\section*{Complications}
\begin{itemize}
    \item Bacterial contamination from poor hygiene, which can make a baby unwell, although this is rare with good practice
    \item Growth of bacteria if milk is left too long at room temperature or in the refrigerator
    \item Loss of some nutritional and immune properties with very long freezer storage
    \item Wasted milk, because it is thawed in too large a quantity or kept too long
    \item Burns to the baby's mouth from milk heated unevenly in a microwave
    \item Freezer burn or an off taste if containers are not sealed well
\end{itemize}

\section*{Prognosis and Outcomes}
With good hygiene and storage practice, expressed breast milk remains safe and nutritious for the recommended storage times, and most babies feed from stored milk without any problem. Breast milk retains many of its benefits even after freezing, and thawed milk is only a little less valuable than fresh. Where guidelines are followed, illness from stored milk is very uncommon, and mothers can confidently build, store, and use a supply of milk that supports their baby's health while allowing flexibility and shared feeding.

\section*{Prevention and Self-Care}
\begin{itemize}
    \item Wash hands and clean all equipment before expressing
    \item Label and date every container at the time of expression
    \item Follow the storage times recommended by your local health service
    \item Store milk at the back of the fridge or freezer and keep the door closed
    \item Thaw milk gently, in the refrigerator or warm water, and never in a microwave
    \item Discard any leftover milk from a bottle after a feed
    \item Ask your midwife, lactation consultant, or maternal and child health nurse if you are ever unsure
\end{itemize}

\section*{Sources and Further Reading}
The following organisations provide reliable, evidence-based information for patients, students, and clinicians.

\begin{itemize}
    \item NHS. \textit{Storing breast milk: guidance on safe storage and handling.} https://www.nhs.uk/conditions/
    \item Centers for Disease Control and Prevention. \textit{Proper storage and preparation of breast milk.} https://www.cdc.gov/
    \item World Health Organization. \textit{Breastfeeding and expressed breast milk guidance.} https://www.who.int/
    \item Mayo Clinic. \textit{Breast milk storage: safe guidelines for keeping expressed milk.} https://www.mayoclinic.org/diseases-conditions/
    \item MedlinePlus. \textit{Storing expressed breast milk: consumer health information.} https://medlineplus.gov/
    \item MSD Manual. \textit{Infant nutrition and the use of expressed breast milk.} https://www.msdmanuals.com/
\end{itemize}
